\section{CourseDocument.cls}\label{sec:ref-coursedocument}

The prose class: syllabus, TA information, policies -- and this manual.
Article-based. Loads Typefaces \texttt{[osf]} by default (the only class
with old-style body figures; \texttt{[lining]} opts out), Titling,
Hyperlinks, InstructorNotes and \file{enumitem} (tunable lists --- a policy
document is more list than prose); no
exercise machinery (a syllabus imports no problems -- add the packages in
the preamble if a prose document ever needs them, as this manual does).
Workflow: section~\ref{sec:prose-documents}.

\subsection*{Class options}

\begin{center}
\begin{tabularx}{\linewidth}{@{}l Y@{}}
\texttt{legalpaper} & legal stock.\\
\texttt{lining} & lining figures in body prose, instead of the class's
  default old-style. For a prose document that is substantially
  \emph{data} -- a syllabus of dates, room numbers, section codes and
  percentages -- where osf reads as clunky. Forwarded to Typefaces
  (section~\ref{sec:ref-typefaces}); it must be a class option because a
  document cannot reach Typefaces, which is loaded during
  \cs{documentclass}.\\
\emph{(other)} & forwarded to \file{article} (\texttt{11pt},
  \texttt{twoside}, \dots).\\
\end{tabularx}
\end{center}

\subsection*{Colours, one name}

The course identity colour is named once -- \env{CourseIdentity}, a
mahogany -- and everything derives from it: \env{HeadingColor} and
\env{TitleColour} (darkened blends), internal link colour, headings.
Restyling a course document is a one-line \cs{colorlet}. The identity is
deliberately not green: green means \emph{solutions / not for students}
across the toolkit (\env{SolutionColor}), and a student-facing document
must not read as green.

\subsection*{Page styles}

Two named styles, neither auto-activated -- \cs{maketitle} runs on
\texttt{plain} so the title page stays bare, then each document opts in:

\begin{center}
\begin{tabularx}{\linewidth}{@{}l Y@{}}
\texttt{toc} & footer page number only; for the contents pages.\\
\texttt{myfancy} & small-caps running section head, footer page number;
  for the body.\\
\end{tabularx}
\end{center}

The canonical opening (used by the syllabus and this manual):
\cs{maketitle}, \cs{pagestyle}\texttt{\{toc\}}, \cs{tableofcontents},
\cs{newpage}, \cs{pagestyle}\texttt{\{myfancy\}}.

\subsection*{Presentation choices}

Sections are numbered three levels deep, set in coloured bold sans via
\file{titlesec}. Paragraphs separate by a blank line, not an indent
(\file{parskip} -- loaded \emph{before} titlesec; the order matters, see
the gotchas in section~\ref{sec:gotchas}). The measure follows the
toolkit-wide rule (\texttt{textwidth} $= 2.6\times$\cs{alphabet},
section~\ref{sec:ref-typefaces}); margins balance 1:1 one-sided and 2:3
two-sided, with a generous marginpar column for instructor notes. The
Titling hooks are set to colour the title and stamp it \emph{at (time)} --
these documents revise mid-term, and the stamp identifies the printing.
